% !PRINTED 410
\weberpage{410}
%<<CONT
und die Summe $\overset{\mathfrak a}{\Sigma}\dfrac{s-1}{(N\mathfrak a)^s}$, über eine Klasse $A$ ausgedehnt, hat für $s=1$ einen bestimmten von Null verschiedenen Grenzwert, der nach § 111, 5. für alle Klassen der gleiche ist. Daraus folgt nach 2.:
\weq{3}{\overset{A}{\Sigma}(\delta,A)=0,}
außer wenn $\delta=1$ oder $=\Delta$ ist,
\weq{4}{\overset{A}{\Sigma}(\delta,A)=h,}
wenn $\delta=1$ oder $=\Delta$ ist [weil $(\Delta,A)=1$, § 104, (10)].

Ferner ist für jede Klasse $A_0$ des Hauptgeschlechts $(\delta,A_0)=1$, während es für jede Klasse $A$, die nicht dem Hauptgeschlecht angehört, wenigstens ein $\delta_1$ gibt, so daß $(\delta_1,A)=-1$ ist. Daraus folgt, wenn $2^\nu$ die Anzahl der Stammteiler ist,
\weq{5}{\begin{aligned}
\overset{\delta}{\Sigma}(\delta,A_0) &= 2^{\nu},\\
(\delta_1,A)\overset{\delta}{\Sigma}(\delta,A) &= \overset{\delta}{\Sigma}(\delta\delta_1,A),
\end{aligned}}
und da $\delta\delta_1$ zugleich mit $\delta$ die Gesamtheit der Stammteiler durchläuft (wenn man $\delta\delta_1$ nach der symbolischen Multiplikation in § 104 reduziert):
\[
(\delta_1,A)\overset{\delta}{\Sigma}(\delta,A) = \overset{\delta}{\Sigma}(\delta,A),
\]
und wenn also $(\delta_1,A)=-1$ ist:
\weq{6}{\overset{\delta}{\Sigma}(\delta,A)=0.}

Hiernach können wir die Doppelsumme
\[
\overset{\delta}{\Sigma}\,\overset{A}{\Sigma}(\delta,A)
\]
auf zwei Arten bestimmen. Es ergibt sich nach (3), (4):
\weq{7}{\overset{\delta}{\Sigma}\,\overset{A}{\Sigma}(\delta,A)=2h,}
und nach (5), (6):
\weq{8}{\overset{\delta}{\Sigma}\,\overset{A}{\Sigma}(\delta,A)=2^{\nu}g,}
wenn $g$ die Anzahl der Klassen des Hauptgeschlechtes (und folglich jedes Geschlechtes) bedeutet. Die Vergleichung gibt:
\weq{9}{h=2^{\nu-1}g,}
also den Satz:

6. \gesperrt{Die Anzahl der Geschlechter ist genau gleich $2^{\nu-1}$, d.\,h. gleich $\tfrac{1}{2}(Z,R)$, § 108, 9.}

\begin{center}\rule{0.32\linewidth}{0.4pt}\end{center}
